\section{I/O管理}

\subsection{I/O设备分类}

\begin{mydef}{I/O设备的分类}{io-device-classification}
计算机的I/O设备按传输速率、信息交换单位、设备共享属性等维度可分为多种类型。

\begin{tablebox}{I/O设备分类}
\centering
\small
\begin{tabularx}{\linewidth}{|p{0.27\linewidth}|p{0.29\linewidth}|X|}
\hline
分类维度 & 类型 & 典型设备 \\
\hline
\multirow{2}{*}{按信息交换单位} & \textbf{块设备（Block Device）} & 磁盘、光盘 \\
                              & \textbf{字符设备（Character Device）} & 键盘、鼠标、打印机 \\
\hline
\multirow{3}{*}{按传输速率} & 低速设备 & 键盘、鼠标等 \\
                            & 中速设备 & 打印机等 \\
                            & 高速设备 & 磁盘、网络接口等 \\
\hline
\multirow{3}{*}{按共享属性} & \textbf{独占设备} & 打印机（临界资源，需互斥使用）\\
                            & \textbf{共享设备} & 磁盘（可多个进程同时交叉访问）\\
                            & \textbf{虚拟设备} & SPOOLing 将独占设备改造为共享设备 \\
\hline
\end{tabularx}
\end{tablebox}

\textbf{块设备 vs 字符设备的核心区别：}
\begin{itemize}
    \item \textbf{块设备}：以数据块为单位传输，支持随机访问（可寻址），有缓冲区缓存机制。如磁盘以扇区（通常512B/4KB）为基本I/O单位。
    \item \textbf{字符设备}：以字符流为单位顺序传输，不提供按数据块随机寻址；系统仍可为其设置字符缓冲或行缓冲。如键盘输入、串口通信。
\end{itemize}
\end{mydef}

\begin{attention}
\textbf{408 辨析：} 块设备 $\neq$ 独占设备。磁盘是块设备同时又是共享设备；打印机是字符设备同时又是独占设备。设备分类维度不同，不可混淆。\textbf{块设备的典型特征是可寻址（可以 lseek），字符设备不可寻址。}
\end{attention}

\subsection{I/O控制方式}

\begin{conclusion}{四种I/O控制方式演进}{io-control-methods}
I/O控制方式的发展本质上是\textbf{将CPU从I/O任务中逐步解放}的过程——每次演进都减少了CPU参与I/O的程度。

\begin{tablebox}{I/O控制方式对比}
\centering
\resizebox{\linewidth}{!}{\begin{tabular}{|l|c|c|c|c|}
\hline
控制方式 & CPU干预频率 & 数据传送单位 & 适用场景 & 特点 \\
\hline
\textbf{程序直接控制（轮询）} & 每字节/每字都要轮询 & 字/字节 & 简单系统 & CPU 忙等，利用率极低 \\
\textbf{中断驱动} & 每字节/每字完成后中断 & 字/字节 & 低速字符设备 & CPU 可并行处理其他任务 \\
\textbf{DMA} & 每块数据传输完成后中断一次 & 块（多字节）& 高速块设备 & CPU 仅参与开始和结束 \\
\textbf{I/O通道} & 每批I/O任务完成后中断一次 & 一组块 & 大型机 & 通道有自己的指令，可执行通道程序 \\
\hline
\end{tabular}}
\end{tablebox}
\end{conclusion}

\subsubsection{程序直接控制（轮询，Polling）}

\begin{mydef}{程序直接控制}{os-programmed-io}
CPU 不断循环检测设备状态寄存器（\texttt{busy/wait} 标志位），当设备就绪时执行一次数据传送（读/写数据寄存器），传送完毕后将设备状态置为就绪（\texttt{command-ready}），然后继续轮询。

\textbf{工作流程：}
\begin{enumerate}
    \item CPU 向设备控制器发送 I/O 命令（如 Read）。
    \item CPU 反复读取状态寄存器，检查设备是否就绪（\texttt{busy == 0}）。
    \item 设备就绪后，CPU 从数据寄存器读取一个字到内存。
    \item 重复第2--3步直到所有数据传输完成。
\end{enumerate}

\textbf{致命缺陷：CPU 忙等（Busy-Waiting）}——在设备完成I/O之前，CPU无法执行任何有意义的计算，CPU利用率极低。
\end{mydef}

\subsubsection{中断驱动方式（Interrupt-Driven I/O）}

\begin{conclusion}{中断驱动I/O}{interrupt-io}
\textbf{核心思想：}CPU 发出 I/O 命令后继续执行其他进程；当设备完成当前数据传输（一个字/字节）后，向 CPU 发出\textbf{中断信号}；CPU 响应中断，在中断处理程序中执行数据传送，然后恢复原进程。

\textbf{工作流程：}
\begin{enumerate}
    \item CPU 向设备控制器发出 Read 命令，然后切换到其他进程运行。
    \item 设备控制器将数据读入自己的数据寄存器，完成后发出中断信号。
    \item CPU 检测到中断，保存当前进程上下文，转去执行中断处理程序。
    \item 中断处理程序将设备数据寄存器中的字/字节传送到内存指定位置。
    \item 若还有数据未传输，CPU 再次启动设备，然后恢复原进程。
\end{enumerate}

\textbf{与轮询的对比：}
\begin{itemize}
    \item 轮询：CPU 主动等设备，浪费 CPU；但响应快（无中断开销）。
    \item 中断：CPU 可以并行做其他事，但每次中断有上下文切换开销。每次只能传输一个字/字节。
\end{itemize}

\textbf{适用场景：}低速字符设备（键盘输入、鼠标、串口通信）。对于高速块设备（磁盘），每次只传一个字且每字都中断，中断过于频繁，开销太大。
\end{conclusion}

\subsubsection{DMA方式（Direct Memory Access）}

\begin{conclusion}{DMA控制方式}{dma-io}
DMA 是\textbf{中断驱动的升级版}——引入专门的 DMA 控制器，在设备和内存之间直接传输\textbf{一整块数据}，CPU 仅在传输开始和结束时参与。数据传输单位从字/字节跃升到块。

\textbf{DMA 控制器的四个寄存器：}
\begin{tablebox}{DMA控制器寄存器}
\centering
\begin{tabularx}{\linewidth}{|p{4.0cm}|X|}
\hline
寄存器 & 功能 \\
\hline
DR（数据寄存器）& 暂存设备与内存之间每次传输的数据 \\
MAR（内存地址寄存器）& 存放目标内存地址，每传输一个字后自动递增 \\
DC（数据计数器）& 存放剩余要传输的字节数/字数，递减至0则传输完成 \\
CR（命令/状态寄存器）& 存放 CPU 发出的命令或设备的状态 \\
\hline
\end{tabularx}
\end{tablebox}

\textbf{DMA 工作流程：}
\begin{enumerate}
    \item CPU 设置 DMA 控制器的 MAR、DC 和命令寄存器，然后继续执行其他进程。
    \item DMA 控制器控制设备将数据块逐字传送到内存：每传一个字，MAR++，DC--。
    \item 当 DC 减至 0 时，DMA 控制器向 CPU 发出\textbf{一次中断}，通知 CPU 传送完成。
    \item CPU 响应中断，执行后处理（如唤醒等待该 I/O 的进程）。
\end{enumerate}
\end{conclusion}

\begin{attention}
\textbf{408 高频辨析：DMA 与中断驱动的关键区别。}
\begin{enumerate}
    \item \textbf{中断频率：}中断驱动每字/字节中断一次；DMA 每块数据仅中断一次。
    \item \textbf{数据传送：}中断驱动由 CPU 执行数据传送（从设备寄存器到内存）；DMA 由 DMA 控制器直接完成，CPU \textbf{不参与数据传输}。
    \item \textbf{总线占用：}DMA 传送期间，DMA 控制器需占用系统总线（周期窃取/周期挪用，Cycle Stealing）。若 CPU 也需要访存，DMA 优先级更高。
    \item \textbf{内存地址：}DMA 传送的是物理地址（不是虚拟地址），因此 DMA 需要连续的物理内存块。
\end{enumerate}
\end{attention}

\begin{formula}{DMA周期窃取时的CPU停顿分析}{dma-cycle-stealing}
设内存周期为 $t_c$，DMA 传送一个字需占用一个内存周期。若 CPU 每执行 $k$ 条指令需访存一次，则 CPU 因 DMA 窃取周期的降速比为：

\[
\text{降速比} = \frac{1}{1 + \frac{1}{k}} \quad \text{（DMA占用总线比例为 } \frac{1}{k+1}\text{）}
\]

当 $k$ 很大（CPU 指令大多访问 Cache）时，DMA 的影响可忽略。这也是现代系统中 Cache 对 DMA 效率的关键意义。
\end{formula}

\subsubsection{I/O通道方式（Channel I/O）}

\begin{mydef}{I/O通道}{io-channel}
\textbf{I/O通道}是一种\textbf{专门负责I/O的协处理器}，有自己的指令集（通道指令），可以执行\textbf{通道程序}来控制一组 I/O 设备完成复杂的 I/O 任务。

\textbf{与 DMA 的区别：}
\begin{itemize}
    \item \textbf{DMA：}只能执行简单的块传输（从设备到固定内存区间），无分支/条件判断能力。
    \item \textbf{通道：}可以执行包含分支、循环的通道程序，能够控制多个设备并发 I/O，CPU 仅需发一条 I/O 指令启动通道即可。
\end{itemize}

\textbf{通道类型：}
\begin{itemize}
    \item \textbf{字节多路通道}：以字节为单位交叉为多台低速设备服务（如多个终端）。
    \item \textbf{数组选择通道}：以块为单位为一台高速设备服务，但通道利用率低（设备独占通道）。
    \item \textbf{数组多路通道}：结合前两者优点，以块为单位交叉为多台高速设备服务（现代大型机主流方案）。
\end{itemize}
\end{mydef}

\begin{attention}
\textbf{408 辨析：} I/O通道方式属于408考试范围，但只在组成原理中详细考查。OS 视角只需了解其层次位置和与DMA的本质区别（通道有指令，可执行通道程序）。
\end{attention}

\subsection{I/O软件层次结构}

\begin{conclusion}{I/O软件的四个层次}{io-software-layers}
I/O 软件采用分层设计，每层对上层屏蔽下层细节，是 OS 分层的典范。

\begin{center}
\begin{tikzpicture}[>=stealth, node distance=1.2cm]
\node[draw, fill=blue!10, minimum width=8cm, minimum height=1cm] (U) {\textbf{用户层 I/O 软件}（\texttt{printf, scanf}等库函数 + SPOOLing 系统）};
\node[draw, fill=green!10, minimum width=8cm, minimum height=1cm, below of=U] (I) {\textbf{设备独立性软件}（命名、保护、阻塞、缓冲、分配）};
\node[draw, fill=yellow!10, minimum width=8cm, minimum height=1cm, below of=I] (D) {\textbf{设备驱动程序}（设备寄存器操作，与硬件直接通信）};
\node[draw, fill=red!10, minimum width=8cm, minimum height=1cm, below of=D] (H) {\textbf{中断处理程序}（保存现场、处理中断、恢复现场）};
\node[draw, fill=gray!10, minimum width=8cm, minimum height=0.8cm, below of=H] (HW) {硬件（设备控制器 + 设备本身）};
\end{tikzpicture}
\end{center}

\noindent\textbf{各层职责详解：}
\begin{enumerate}
    \item \textbf{中断处理程序（Interrupt Handler）——最底层软件}
    \begin{itemize}
        \item I/O 完成后，设备控制器发出中断信号。
        \item CPU 响应中断：硬件保存 PC 和 PSW，然后跳转到中断向量表指定的中断处理程序入口。
        \item 中断处理程序：保存寄存器上下文 $\to$ 处理中断（读取设备状态，唤醒等待进程）$\to$ 恢复上下文 $\to$ 中断返回。
        \item \textbf{关键要求：中断处理程序必须尽可能短}，复杂性推迟到下半部（Bottom Half）处理。
    \end{itemize}

    \item \textbf{设备驱动程序（Device Driver）——设备相关}
    \begin{itemize}
        \item 每种设备（或每类设备）对应一个驱动程序，由设备制造商提供。
        \item 负责将 OS 的通用 I/O 请求（如 ``读第 n 块''）翻译为设备控制器的具体命令序列（操作设备寄存器）。
        \item 向上提供统一接口（open, read, write, ioctl, release），向下操作设备控制器的\textbf{控制寄存器、状态寄存器、数据寄存器}。
        \item \textbf{驱动程序是内核的一部分}（运行在内核态），但代码量和 bug 占比极高。
    \end{itemize}

    \item \textbf{设备独立性软件（Device-Independent Software）——核心层}
    \begin{itemize}
        \item 执行所有设备共用的 I/O 功能，为上层提供\textbf{统一的逻辑设备名}到物理设备的映射。
        \item \textbf{主要功能：}
        \begin{itemize}
            \item \textbf{设备命名与保护}：将逻辑设备名（如 \texttt{/dev/sda}）映射到对应的驱动程序。
            \item \textbf{设备分配与回收}：管理独占设备（如打印机）的分配，防止死锁。
            \item \textbf{缓冲管理}：实现单缓冲、双缓冲、缓冲池，平滑 CPU 和 I/O 设备的速度差异。
            \item \textbf{差错处理}：处理一般的 I/O 错误（如磁盘读错重试）。
            \item \textbf{提供统一块大小}：向上层屏蔽不同磁盘扇区大小的差异。
        \end{itemize}
        \item \textbf{典型实现：}UNIX/Linux 的文件系统层即承担了设备独立性软件的角色。
    \end{itemize}

    \item \textbf{用户层 I/O 软件（User-Level I/O Software）}
    \begin{itemize}
        \item 库函数：\texttt{printf} 调用 \texttt{write} 系统调用，\texttt{scanf} 调用 \texttt{read}。
        \item SPOOLing 系统：在用户层将独占设备虚拟化为共享设备（详见下一节）。
        \item 假脱机管理程序（如打印守护进程 \texttt{lpd}）通常以用户进程形式运行。
    \end{itemize}
\end{enumerate}
\end{conclusion}

\begin{attention}
\textbf{408 高频辨析：各层的归属。}
\begin{itemize}
    \item 中断处理程序、设备驱动程序 \textbf{运行在内核态}（属于内核）。
    \item 设备独立性软件 \textbf{运行在内核态}（属于内核）。
    \item SPOOLing 系统 \textbf{通常在用户层}（以守护进程形式），但也涉及内核态的系统调用。
\end{itemize}

\textbf{关键题眼：} ``向上提供统一接口，向下操作设备寄存器'' $\to$ 这是在描述 \textbf{设备驱动程序}。
\end{attention}

\subsection{SPOOLing系统（假脱机技术）}

\begin{mydef}{SPOOLing}{spooling-def}
\textbf{SPOOLing（Simultaneous Peripheral Operations On-Line，假脱机/外部设备联机并行操作）} 是一种将\textbf{独占设备（如打印机）改造为可共享的虚拟设备}的技术。

\textbf{核心思想：}用磁盘（共享设备）作为中转站，截获对独占设备的输出请求，先写入磁盘上的\textbf{SPOOLing 文件}，再由 SPOOLing 守护进程在后台逐个传输到物理设备。用户进程以为自己在独占设备，实际只是写入了磁盘。

\begin{center}
\begin{tikzpicture}[>=stealth, node distance=2.5cm]
\node[draw, fill=yellow!20, minimum width=3cm] (P1) {进程 P1};
\node[draw, fill=yellow!20, minimum width=3cm, below of=P1, yshift=0.5cm] (P2) {进程 P2};
\node[draw, fill=gray!20, minimum width=4cm, right of=P1, xshift=2cm] (DISK) {磁盘（SPOOLing 井）};
\node[draw, fill=green!20, minimum width=4cm, right of=DISK, xshift=2cm] (DAEMON) {SPOOLing 守护进程};
\node[draw, fill=red!20, minimum width=3cm, right of=DAEMON, xshift=1.5cm] (DEV) {打印机};
\draw[->] (P1) -- node[above] {输出到井} (DISK);
\draw[->] (P2) -- node[below] {输出到井} (DISK);
\draw[->] (DISK) -- node[above] {逐个取出} (DAEMON);
\draw[->] (DAEMON) -- node[above] {实际打印} (DEV);
\end{tikzpicture}
\end{center}
\end{mydef}

\begin{conclusion}{SPOOLing系统的组成}{spooling-components}
SPOOLing 系统由三部分组成：

\begin{enumerate}
    \item \textbf{输入井和输出井}：磁盘上开辟的两块区域。
    \begin{itemize}
        \item \textbf{输入井（Input Well）}：暂存输入设备（如读卡机）的数据，供进程按需读取。
        \item \textbf{输出井（Output Well）}：暂存进程的输出数据，等待输出设备空闲时送出。
    \end{itemize}

    \item \textbf{输入进程和输出进程（SPOOLing守护进程）}：
    \begin{itemize}
        \item \textbf{输入进程（预输入程序/SPi）}：将数据从物理输入设备预先读入输入井。
        \item \textbf{输出进程（缓输出程序/SPo）}：当输出设备空闲时，将输出井中的数据送到物理设备。
        \item 输出进程通常以后台守护进程形式运行。
    \end{itemize}

    \item \textbf{井管理程序}：管理输入/输出井中的空间分配和数据调度。
\end{enumerate}
\end{conclusion}

\begin{conclusion}{SPOOLing的工作原理与效果}{spooling-workflow}
\textbf{以共享打印机为例的完整流程：}

\begin{enumerate}
    \item 用户进程调用 \texttt{print()} —— 系统并不直接分配打印机，而是在输出井中创建 SPOOLing 文件，将打印数据写入磁盘。
    \item 用户进程感觉 ``打印完成''，继续运行（实际上只是写盘完成）。
    \item 打印守护进程（如 UNIX 的 \texttt{lpd}）轮询输出井，发现有待打印文件，且打印机能接收数据时，将文件内容逐个发送给打印机。
    \item 多个进程的打印请求在磁盘上排队（Spooling Queue），物理打印机串行服务。
\end{enumerate}

\textbf{SPOOLing 的效果：}
\begin{itemize}
    \item \textbf{将独占设备改造为共享设备}：多个进程可同时 ``使用'' 打印机（实际是同时写磁盘）。
    \item \textbf{将低速 I/O 操作变为高速 I/O 操作}：对于用户进程，写磁盘比等待打印机快得多。
    \item \textbf{实现了虚拟设备功能}：每个进程都觉得自己独占了打印机。
\end{itemize}
\end{conclusion}

\begin{attention}
\textbf{408 核心考点：SPOOLing 的本质。}
\begin{enumerate}
    \item SPOOLing 的``假脱机''中，``假''在\textbf{用户进程以为在与物理 I/O 设备直接交互，实际上是在与磁盘交互}。
    \item SPOOLing 输入/输出井 \textbf{位于磁盘上}（不是内存中！）。
    \item SPOOLing 是一种典型的\textbf{虚拟化技术}——与虚拟内存（将磁盘空间当内存用）互为镜像：\textbf{虚拟内存把磁盘当内存用；SPOOLing 把磁盘当设备用。}
    \item \textbf{SPOOLing 的常见辨析：}其典型用途是把打印机等独占设备改造为可由多个进程并发申请的虚拟共享设备；磁盘本身支持按块交错访问，通常不需要靠 SPOOLing 才能实现共享。
\end{enumerate}
\end{attention}

\subsection{缓冲区管理}

\begin{mydef}{缓冲的引入动机}{buffer-motivation}
\textbf{为什么需要缓冲？}
\begin{enumerate}
    \item \textbf{缓和 CPU 与 I/O 设备的速度不匹配}：CPU 处理速度远快于 I/O 设备，缓冲让两者可以并行工作。
    \item \textbf{减少对 CPU 的中断频率}：没有缓冲时，每传一个字节就中断一次；有缓冲后，可等缓冲区满（或半满）再中断。
    \item \textbf{解决 DMA 和中断的数据粒度不匹配}：DMA 按块传输，但进程可能按字节处理数据，缓冲区桥接了两者。
    \item \textbf{提高 CPU 和 I/O 设备的并行性}：CPU 处理上一块数据的同时，设备可以将下一块数据送入缓冲区。
\end{enumerate}

\textbf{缓冲实现的位置：}缓冲区通常在内核空间的内存中分配。
\end{mydef}

\subsubsection{单缓冲（Single Buffer）}

\begin{conclusion}{单缓冲的工作原理与时间分析}{single-buffer}
\textbf{结构：}在设备与用户进程之间设置一个缓冲区。设备和进程轮流使用该缓冲区——两者\textbf{不能同时操作同一缓冲区}。

\textbf{设以下参数：}
\begin{itemize}
    \item $T$：设备将一块数据输入到缓冲区所需时间（I/O传输时间）
    \item $M$：OS 将缓冲区的数据传送到用户工作区所需时间（数据搬移时间）
    \item $C$：CPU/用户进程处理一块数据所需时间
\end{itemize}

\textbf{单缓冲下的数据处理时序（以输入为例）：}

\begin{center}
\begin{tikzpicture}[>=stealth]
\draw[->] (0,0) -- (12,0) node[right] {$t$};
% Block 1
\draw[fill=blue!20] (0,0.5) rectangle (2,1.2) node[midway] {\small $T$ (读块1)};
\draw[fill=green!20] (2,0.5) rectangle (2.8,1.2) node[midway] {\tiny $M$};
\draw[fill=yellow!20] (2.8,0.5) rectangle (5.2,1.2) node[midway] {\small $C$ (处理块1)};
% Block 2
\draw[fill=blue!20] (5.2,-0.5) rectangle (7.2,0.2) node[midway] {\small $T$ (读块2)};
\draw[fill=green!20] (7.2,-0.5) rectangle (8,0.2) node[midway] {\tiny $M$};
\draw[fill=yellow!20] (8,-0.5) rectangle (10.4,0.2) node[midway] {\small $C$ (处理块2)};
\end{tikzpicture}
\end{center}

\text{相邻两块完成处理的间隔（稳态）}：
\[
\text{每块时间} = \max(C, T) + M
\]

\textbf{原因分析：}
\begin{itemize}
    \item 把当前块从系统缓冲区搬到用户区需 $M$，搬完后缓冲区才可供设备输入下一块。
    \item 下一块的输入 $T$ 可与当前块的处理 $C$ 重叠，二者的较大值决定等待时间。
    \item 搬移阶段 $M$ 使用唯一的系统缓冲区，不能与设备向该缓冲区输入下一块重叠。
\end{itemize}

\textbf{处理 $N$ 块数据的总时间：}
\[
T_{\text{single}} = T + M + C +(N-1)[\max(C,T)+M]
\]
第一块先经历 $T+M+C$；之后每增加一块，完成时刻增加 $\max(C,T)+M$。
\end{conclusion}

\subsubsection{双缓冲（Double Buffer）}

\begin{conclusion}{双缓冲的工作原理与时间分析}{double-buffer}
\textbf{结构：}设置两个缓冲区：Buffer0 和 Buffer1（又称\textbf{缓冲对换}，Buffer Swapping）。设备向一个缓冲区填数据时，CPU 可从另一个缓冲区取数据并处理。

\textbf{双缓冲下的数据处理时序：}

\begin{center}
\begin{tikzpicture}[>=stealth]
\draw[->] (0,0) -- (12,0) node[right] {$t$};
% Buffer0: device fills, CPU processes
\draw[fill=blue!20] (0,0.5) rectangle (2,1.2) node[midway] {\small $T$ (B0: 读块1)};
\draw[fill=green!20] (2,0.5) rectangle (2.8,1.2) node[midway] {\tiny $M$};
\draw[fill=yellow!20] (2.8,0.5) rectangle (5.2,1.2) node[midway] {\small $C$ (B0: 处理块1)};
% Buffer1: device fills in parallel
\draw[fill=blue!20] (2,-0.5) rectangle (4,0.2) node[midway] {\small $T$ (B1: 读块2)};
\draw[fill=green!20] (4,-0.5) rectangle (4.8,0.2) node[midway] {\tiny $M$};
\draw[fill=yellow!20] (5.2,-0.5) rectangle (7.6,0.2) node[midway] {\small $C$ (B1: 处理块2)};
\end{tikzpicture}
\end{center}

\text{相邻两块完成处理的间隔（稳态）}：
\[
\text{每块时间} = \max(T, M+C)
\]

\textbf{原因分析：}
\begin{itemize}
    \item 设备可以向一个缓冲区输入（耗时 $T$），同时系统从另一个缓冲区搬移并由进程处理上一块（串行耗时 $M+C$）。
    \item 两条并行流水路径中较慢的一条决定稳态吞吐率，因此比较的是 $T$ 与 $M+C$，而不是只比较 $T$ 与 $C$。
\end{itemize}

\textbf{处理 $N$ 块数据的总时间：}
\[
T_{\text{double}} = T + M + C +(N-1)\max(T,M+C)
\]
第一块仍需完整的 $T+M+C$；从第二块起，设备输入和“搬移+处理”两条路径重叠。
\end{conclusion}

\begin{attention}
\textbf{408 计算题关键区分——单缓冲 vs 双缓冲：}

\noindent\fcolorbox{red}{white}{\begin{minipage}{0.9\linewidth}
\textbf{真题场景：} 文件占 $N$ 个磁盘块，每块大小等于缓冲区大小。分别计算单缓冲和双缓冲下读入并处理整个文件的时间。

\textbf{单缓冲：} $T_{\text{single}}=T+M+C+(N-1)[\max(C,T)+M]$

\textbf{双缓冲：} $T_{\text{double}}=T+M+C+(N-1)\max(T,M+C)$

\textbf{陷阱：}双缓冲的 CPU 路径是 $M+C$，两者通常不能彼此重叠；不能把公式误写为 $\max(T,M,C)$。若题目未给 $M$，应按题干明确的假设取 $M=0$ 或认为它已包含在其他时间中。
\end{minipage}}
\end{attention}

\subsubsection{循环缓冲（Circular Buffer）}

\begin{mydef}{循环缓冲}{circular-buffer}
\textbf{结构：}在双缓冲的基础上，将多个缓冲区组织成一个环形队列。每个缓冲区有一个状态标记：空（Empty）、满（Full）、正在使用（In-Use）。

\begin{center}
\begin{tikzpicture}
\foreach \i in {0,...,7} {
    \node[draw, circle, minimum size=1cm] (B\i) at ({360/8*\i}:2cm) {Buf\i};
}
\draw[->] (B0) -- (B1);
\draw[->] (B1) -- (B2);
\draw[->] (B2) -- (B3);
\draw[->] (B3) -- (B4);
\draw[->] (B4) -- (B5);
\draw[->] (B5) -- (B6);
\draw[->] (B6) -- (B7);
\draw[->] (B7) -- (B0);
\node[draw, fill=yellow!20, below=1.5cm] at (B3) {in: 设备写入指针};
\node[draw, fill=green!20, below=1.5cm] at (B7) {out: CPU读出指针};
\draw[->, thick, blue] (B3) -- node[right] {\small \texttt{in}} ++(0,-1);
\draw[->, thick, red] (B7) -- node[right] {\small \texttt{out}} ++(0,-1);
\end{tikzpicture}
\end{center}

\textbf{关键指针：}
\begin{itemize}
    \item \textbf{NextIn（输入指针）}：指向设备下一个要填入数据的空缓冲区。
    \item \textbf{NextOut（输出指针）}：指向 CPU 下一个要取数据的满缓冲区。
    \item \textbf{Current}：指向当前正在被 CPU 处理的缓冲区（已从 NextOut 取出）。
\end{itemize}

\textbf{与生产者-消费者的关系：}循环缓冲本质上是生产者-消费者问题在 I/O 缓冲中的具体应用。设备是生产者（填满缓冲区），CPU 是消费者（取空缓冲区）。管理和同步由 OS 的缓冲区管理程序负责。
\end{mydef}

\subsubsection{缓冲池（Buffer Pool）}

\begin{conclusion}{缓冲池}{buffer-pool}
\textbf{缓冲池}是系统级的大型缓冲区集合，可以被多种设备和进程共享使用，而不是专属于某个特定设备。

\textbf{缓冲池中的三种缓冲区队列：}
\begin{tablebox}{缓冲池的三种队列}
\centering
\begin{tabularx}{\linewidth}{|p{0.25\linewidth}|p{0.22\linewidth}|X|}
\hline
队列类型 & 状态 & 用途 \\
\hline
空缓冲队列（emq）& 空闲，可分配 & 设备和进程均可从此获取空缓冲区 \\
输入队列（inq）& 装满输入数据 & 设备将数据填入后挂入此队列；CPU从此取出数据 \\
输出队列（outq）& 装满输出数据 & 进程将输出数据填入后挂入此队列；设备从此取出并输出 \\
\hline
\end{tabularx}
\end{tablebox}

\textbf{缓冲池的操作原语：}
\begin{itemize}
    \item \textbf{收容输入（GetBuf\_Input）}：从 emq 取一个空缓冲区 $\to$ 挂入 inq。
    \item \textbf{提取输入（ReleaseBuf\_Input）}：从 inq 取一个满缓冲区 $\to$ 供进程使用 $\to$ 回收后挂入 emq。
    \item \textbf{收容输出（GetBuf\_Output）}：从 emq 取一个空缓冲区 $\to$ 挂入 outq。
    \item \textbf{提取输出（ReleaseBuf\_Output）}：从 outq 取一个满缓冲区 $\to$ 供设备输出 $\to$ 回收后挂入 emq。
\end{itemize}

\textbf{与专用缓冲的对比：}
\begin{itemize}
    \item \textbf{专用缓冲（单/双/循环缓冲）}：为特定设备静态分配，利用率可能不足（设备空闲时缓冲区浪费）或不足（高峰期不够用）。
    \item \textbf{缓冲池}：所有设备共享，动态分配，利用率高，但管理开销更大。
\end{itemize}
\end{conclusion}

\begin{attention}
\textbf{408 辨析：四个缓冲区概念的区别。}
\begin{itemize}
    \item \textbf{单缓冲}：稳态间隔为 $\max(C,T)+M$。
    \item \textbf{双缓冲}：稳态间隔为 $\max(T,M+C)$，适合连续 I/O（如磁盘顺序读）。
    \item \textbf{循环缓冲}：多个缓冲区环形组织，适合生产者-消费者速率不稳定的场景。
    \item \textbf{缓冲池}：系统全局共享的缓冲区集合，通过收容/提取操作管理。
\end{itemize}
\textbf{408 经常考察：} 给定 $C$ 和 $T$ 的值，判断单缓冲和双缓冲的总处理时间差异。当 $C$ 和 $T$ 接近时，双缓冲的优势最大（并行度最高）。
\end{attention}

\subsection{408 高频考点速查}

\begin{conclusion}{I/O管理 — 408常考点}{os5-keypoints}
\begin{enumerate}
    \item \textbf{四种 I/O 控制方式}：轮询$\to$中断$\to$DMA$\to$通道，每次演进减少了 CPU 参与 I/O 的程度和中断频率。DMA 与中断的最关键区别：\textbf{DMA 的数据传输不由 CPU 完成，且每块只中断一次}。

    \item \textbf{I/O 软件层次}：四层自底向上依次是中断处理程序$\to$设备驱动程序$\to$设备独立性软件$\to$用户层 I/O 软件。SPOOLing 属于用户层（守护进程形式）。

    \item \textbf{SPOOLing 系统}：独占设备$\to$共享设备的虚拟化技术。输入/输出井\textbf{在磁盘上}。原理：截获输出请求写入磁盘井，后台守护进程逐个送往物理设备。

    \item \textbf{缓冲区处理时间计算}：
    \begin{itemize}
        \item 单缓冲：$T+M+C+(N-1)[\max(C,T)+M]$
        \item 双缓冲：$T+M+C+(N-1)\max(T,M+C)$
    \end{itemize}

    \item \textbf{块设备 vs 字符设备}：块设备可寻址（支持 lseek），以块为单位；字符设备不可寻址，以字符流为单位。

    \item \textbf{DMA 周期窃取}：DMA 传送期间需占用总线，与 CPU 争用内存周期，但每块只窃取少量周期，总体影响不大。
\end{enumerate}
\end{conclusion}

\newpage
\section{习题}

\subsection{基础过关题}

\begin{enumerate}

\item \textbf{（真题风格·选择题）} 系统将数据从磁盘读到内存的过程包括以下操作：
\begin{center}
\begin{tabular}{l}
（1） DMA 控制器发出中断请求 \\
（2） 初始化 DMA 控制器并启动磁盘 \\
（3） 从磁盘传输一块数据到内存缓冲区 \\
（4） 执行「DMA 结束」中断服务程序
\end{tabular}
\end{center}
正确的执行顺序是
\begin{center}
\begin{tabular}{ll}
(A) （3）$\to$（1）$\to$（2）$\to$（4） & (B) （2）$\to$（3）$\to$（1）$\to$（4） \\[6pt]
(C) （2）$\to$（1）$\to$（3）$\to$（4） & (D) （1）$\to$（2）$\to$（4）$\to$（3）
\end{tabular}
\end{center}

\item \textbf{（真题风格·选择题）} 在系统内存中设置磁盘缓冲区的主要目的是
\begin{center}
\begin{tabular}{ll}
(A) 减少磁盘 I/O 次数 & (B) 减少平均寻道时间 \\[6pt]
(C) 提高磁盘数据可靠性 & (D) 实现设备无关性
\end{tabular}
\end{center}

\item \textbf{（真题风格·选择题）} 下列关于驱动程序的叙述中，\textbf{不正确}的是
\begin{center}
\begin{tabular}{ll}
(A) 驱动程序与 I/O 控制方式无关 \\[4pt]
(B) 初始化设备是由驱动程序控制的 \\[4pt]
(C) 驱动程序执行 I/O 操作时，进程可能会阻塞 \\[4pt]
(D) 对磁盘的读/写操作最终转换为驱动程序中的读/写函数
\end{tabular}
\end{center}

\item \textbf{（填空题）} I/O 软件的四个层次自底向上依次为：\underline{\qquad}、\underline{\qquad}、\underline{\qquad} 和 \underline{\qquad}。其中，SPOOLing 系统属于 \underline{\qquad} 层；设备独立性软件负责将 \underline{\qquad} 设备名映射到物理设备；向上提供统一接口、向下操作设备寄存器的是 \underline{\qquad}。

\item \textbf{（简答题）} 简述 DMA 控制方式与中断驱动方式的核心区别（至少三点），并说明为什么 DMA 方式更适合高速块设备。

\end{enumerate}

\subsection{真题风格与综合训练}

\begin{enumerate}[resume]

\item \textbf{（真题风格·选择题）} 下列关于 SPOOLing 技术的叙述中，\textbf{错误}的是
\begin{center}
\begin{tabular}{ll}
(A) SPOOLing 技术需要外存的支持 \\[4pt]
(B) SPOOLing 技术需要多道程序技术的支持 \\[4pt]
(C) SPOOLing 技术可以使多个作业共享一台独占设备 \\[4pt]
(D) SPOOLing 技术由用户作业控制设备与输入/输出井之间的数据传送
\end{tabular}
\end{center}

\item \textbf{（真题风格·选择题）} 用户程序发出磁盘 I/O 请求后，系统的正确处理流程是
\begin{center}
\begin{tabular}{ll}
(A) 用户程序$\to$系统调用处理程序$\to$设备驱动程序$\to$中断处理程序 \\[4pt]
(B) 用户程序$\to$系统调用处理程序$\to$中断处理程序$\to$设备驱动程序 \\[4pt]
(C) 用户程序$\to$设备驱动程序$\to$系统调用处理程序$\to$中断处理程序 \\[4pt]
(D) 用户程序$\to$设备驱动程序$\to$中断处理程序$\to$系统调用处理程序
\end{tabular}
\end{center}

\item \textbf{（真题风格·选择题）} 设系统缓冲区和用户工作区均采用单缓冲，从外设读入 1 个数据块到系统缓冲区的时间为 100，从系统缓冲区读入 1 个数据块到用户工作区的时间为 5，对用户工作区中的 1 个数据块进行分析的时间为 90。进程从外设读入并分析 2 个数据块的最短时间是
\begin{center}
\begin{tabular}{ll}
(A) 200 & (B) 295 \\[6pt]
(C) 300 & (D) 390
\end{tabular}
\end{center}

\item \textbf{（真题风格·选择题）} 某文件占 10 个磁盘块，现要把该文件磁盘块逐个读入主存缓冲区，并送用户区进行分析。假设一个缓冲区与一个磁盘块大小相同，把一个磁盘块读入缓冲区的时间为 100 $\mu\text{s}$，将缓冲区的数据传送到用户区的时间是 50 $\mu\text{s}$，CPU 对一块数据进行分析的时间为 50 $\mu\text{s}$。在单缓冲区和双缓冲区结构下，读入并分析完该文件的时间分别是
\begin{center}
\begin{tabular}{ll}
(A) 1500 $\mu\text{s}$、1000 $\mu\text{s}$ & (B) 1550 $\mu\text{s}$、1100 $\mu\text{s}$ \\[6pt]
(C) 1550 $\mu\text{s}$、1550 $\mu\text{s}$ & (D) 2000 $\mu\text{s}$、2000 $\mu\text{s}$
\end{tabular}
\end{center}

\item \textbf{（真题风格，综合题）} 某文件占 200 个磁盘块，磁盘每块大小为 1 KB，缓冲区大小等于磁盘块大小。已知：磁盘将一块数据读入缓冲区的时间 $T = 0.2$ ms，系统将缓冲区数据传送到用户区的时间 $M = 0.01$ ms，CPU 处理一块数据的时间 $C = 0.3$ ms。

\begin{enumerate}
    \item[(1)] 分别计算使用单缓冲和双缓冲时，读入并处理完整个文件需要的总时间。
    \item[(2)] 若缓冲区大小为 2 KB（= 2 个磁盘块），在双缓冲下，读入并处理完整个文件需要多长时间？（设 $T$ 与块大小成正比，CPU 处理时间 $C$ 不变。）
    \item[(3)] 以上两种缓冲大小下，哪种配置的 I/O 吞吐率更高？为什么？
\end{enumerate}

\end{enumerate}

\vspace{8pt}
\noindent\textbf{解题提示：}
\begin{itemize}
    \item 选择题中的单/双缓冲计算使用：单缓冲 $T+M+C+(N-1)[\max(C,T)+M]$；双缓冲 $T+M+C+(N-1)\max(T,M+C)$。
    \item 随后的综合题中，$\max(T,M+C)=\max(0.2,0.31)=0.31$ ms，双缓冲总时间为 $0.2+0.01+0.3+199\times0.31=62.20$ ms。
\end{itemize}

\vspace{8pt}
\noindent\textbf{拓展阅读：}
\begin{itemize}
    \item OSTEP 第 36--37 章「I/O设备」和「磁盘驱动器」——从设备控制器、驱动程序到 RAID 的完整故事。
    \item 陈海波《现代操作系统：原理与实现》第 8 章「设备管理」——Linux 设备驱动模型和中断下半部机制。
    \item Brooks《人月神话》第 2 章（名著阅读）——经典的设备独立性思想来源。
\end{itemize}
